% ============================================================
%  CODE APPENDIX
%
%  main.tex resets the listing counter and gives it an "S" prefix
%  just before including this file, so the listings here are S1,
%  S2, ... and do not clash with the excerpts in the chapters.
%
%  The listing style (colours, line numbers, frame, line breaking)
%  is set once in main.tex as "report_listing" — change it there,
%  not here.
% ============================================================

\section*{Python code}
\phantomsection
\addcontentsline{toc}{section}{Python code}

The complete program used for Problems 3 and 4 is listed below. It is written
to be run as a single file: the two \texttt{main\_*} functions at the bottom
produce the figures referenced in those chapters. Replace it with your own
source — the listing environment breaks long lines automatically and marks each
break with an arrow, so nothing runs off the page.

\begin{lstlisting}[language=Python, caption=Newton basins and the Rossler attractor, label=lst:p5_full]
"""Basins of attraction for Newton's method, and a Rossler trajectory.

Run directly:  python report_figures.py
Writes two PNG files into the working directory.
"""

import numpy as np
import matplotlib
matplotlib.use('Agg')          # no GUI window; write files only
from matplotlib import pyplot as plt

############################################################
# Problem 3: Newton's method on p(z) = z^3 - 1
############################################################

TOLERANCE      = 1e-8          # convergence test on |z_{n+1} - z_n|
MIN_DERIVATIVE = 1e-12         # guard against division by zero
MAX_ITERATIONS = 50

# The three cube roots of unity, exp(2*pi*i*k/3) for k = 0, 1, 2.
ROOTS = np.exp(2j * np.pi * np.arange(3) / 3)


def newton_basins(resolution=1200, extent=1.75):
    """Return (root_index, iteration_count) for a grid of starting points.

    root_index is -1 wherever the iteration failed to converge, which
    happens on the fractal boundary and at the origin, where p'(z) = 0.
    """
    axis = np.linspace(-extent, extent, resolution)
    real, imaginary = np.meshgrid(axis, axis)
    z = real + 1j * imaginary

    root_index = np.full(z.shape, -1, dtype=int)
    iterations = np.zeros(z.shape, dtype=int)
    active     = np.ones(z.shape, dtype=bool)

    for step in range(MAX_ITERATIONS):
        derivative = 3.0 * z[active] ** 2

        # Points where the derivative vanishes can never be stepped.
        stuck = np.abs(derivative) < MIN_DERIVATIVE
        if stuck.any():
            still_active        = np.flatnonzero(active)
            active.flat[still_active[stuck]] = False
            derivative          = derivative[~stuck]

        previous   = z[active]
        z[active]  = previous - (previous ** 3 - 1.0) / derivative
        converged  = np.abs(z[active] - previous) < TOLERANCE

        if converged.any():
            index = np.flatnonzero(active)[converged]
            # Assign each converged point to the nearest cube root.
            distance = np.abs(z.flat[index][:, None] - ROOTS[None, :])
            root_index.flat[index] = np.argmin(distance, axis=1)
            iterations.flat[index] = step + 1
            active.flat[index]     = False

        if not active.any():
            break

    return root_index, iterations


def plot_basins(root_index, iterations, filename='newton_fractal.png'):
    """Colour by root, shade by how many iterations that root took."""
    shade = 1.0 - iterations / max(iterations.max(), 1)
    image = np.zeros(root_index.shape + (3,))
    for channel in range(3):
        image[..., channel] = np.where(root_index == channel, shade, 0.0)

    figure, axes = plt.subplots(figsize=(7, 7))
    axes.imshow(image, extent=[-1.75, 1.75, -1.75, 1.75], origin='lower')
    axes.set_xlabel('Re(z)')
    axes.set_ylabel('Im(z)')
    axes.set_title('Newton fractal - z^3 - 1 = 0')
    figure.tight_layout()
    figure.savefig(filename, dpi=300)
    plt.close(figure)

############################################################
# Problem 4: the Rossler system, integrated with RK4
############################################################

A_PARAMETER = 0.2
B_PARAMETER = 0.2
C_PARAMETER = 5.7


def rossler_derivative(state):
    """The right-hand side of the three coupled equations."""
    x, y, z = state
    return np.array([
        -y - z,
        x + A_PARAMETER * y,
        B_PARAMETER + z * (x - C_PARAMETER),
    ])


def runge_kutta_4(state, delta_time):
    """One classical fourth-order Runge-Kutta step."""
    k1 = rossler_derivative(state)
    k2 = rossler_derivative(state + 0.5 * delta_time * k1)
    k3 = rossler_derivative(state + 0.5 * delta_time * k2)
    k4 = rossler_derivative(state + delta_time * k3)
    return state + (delta_time / 6.0) * (k1 + 2 * k2 + 2 * k3 + k4)


def integrate_rossler(total_time=500.0, delta_time=0.005, transient=0.10):
    """Integrate one trajectory and drop the initial transient."""
    n_steps    = int(total_time / delta_time)
    trajectory = np.zeros((n_steps, 3))
    state      = np.array([1.0, 1.0, 1.0])

    for step in range(n_steps):
        state             = runge_kutta_4(state, delta_time)
        trajectory[step]  = state

    return trajectory[int(transient * n_steps):]


def fixed_points():
    """The two real fixed points, from the quadratic x^2 - cx + ab = 0."""
    discriminant = C_PARAMETER ** 2 - 4 * A_PARAMETER * B_PARAMETER
    if discriminant < 0:
        return []                      # no real fixed point for these parameters

    points = []
    for sign in (+1, -1):
        x = (C_PARAMETER + sign * np.sqrt(discriminant)) / 2
        points.append((x, -x / A_PARAMETER, x / A_PARAMETER))
    return points


def plot_rossler(trajectory, filename='rossler.png'):
    """Plot the trajectory in 3D, coloured by time."""
    figure = plt.figure(figsize=(7, 7))
    axes   = figure.add_subplot(111, projection='3d')

    # One coloured segment per step: matplotlib has no per-vertex colour
    # for 3D lines, so the trajectory is drawn in chunks.
    chunk = 200
    colours = plt.cm.turbo(np.linspace(0, 1, len(trajectory) // chunk + 1))
    for index, start in enumerate(range(0, len(trajectory) - 1, chunk)):
        piece = trajectory[start:start + chunk + 1]
        axes.plot(piece[:, 0], piece[:, 1], piece[:, 2],
                  color=colours[index], linewidth=0.6)

    axes.set_xlabel('x')
    axes.set_ylabel('y')
    axes.set_zlabel('z')
    axes.set_title('The Rossler Attractor')
    figure.tight_layout()
    figure.savefig(filename, dpi=300)
    plt.close(figure)

############################################################
# Entry points
############################################################

def main_problem_3():
    root_index, iterations = newton_basins()
    plot_basins(root_index, iterations)
    unassigned = np.count_nonzero(root_index < 0)
    print(f'Unassigned points: {unassigned} '
          f'({100 * unassigned / root_index.size:.2f} percent of the grid)')


def main_problem_4():
    trajectory = integrate_rossler()
    plot_rossler(trajectory)
    for point in fixed_points():
        print('Fixed point: ({:.4f}, {:.4f}, {:.4f})'.format(*point))


if __name__ == '__main__':
    main_problem_3()
    main_problem_4()
\end{lstlisting}
